% !PRINTED 265
\weberpage{265}
%<<CONT
in which the $M_1, M_2 \ldots M_k$ depend rationally on the $a_{i,k}, b_{i,k}$. In particular,
\[
M_k = \Sigma \pm b_{1,1} b_{2,2} \ldots b_{k,k},
\]
and one can always choose the $b_{i,k}$ so that $M_k$ is different from zero. By \S.\,31 one can therefore choose $\varepsilon$ so that, if $R_k$ is different from zero, $R_k'$ has the same sign as $R_k$, and if $R_k$ vanishes, $R_k'$ does not vanish.

We can combine what has thus been proved with the results of \S.\,81 into a general rule for determining the number of negative squares, valid also in the case of vanishing determinant.

% The source continues to use the same variant rho (varrho) introduced in section 80.
VI. \gesperrt{To determine the numbers $\pi$, $\nu$, $\varrho$ of the function $\varphi(x)$, arrange the variables $x_1, x_2 \ldots x_m$ so that in the chain of principal subdeterminants}
\weq{7}{R_m, R_{m-1} \ldots R_1, R_0}
\gesperrt{as small a number of initial terms as possible vanishes, and so that of the following terms no two standing next to one another vanish; $\varrho$ is then the number of vanishing initial terms, $\nu$ the number of changes of sign, and $\pi = m - \nu - \varrho$.}

Finally let it be remarked that the number of changes of sign in the series (7) can also be counted from right to left, that is, one finds the same number of changes of sign if the series is written in the reverse order.

\begin{center}
\S.\,84.

\gesperrt{Application to the Bezoutiant.}
\end{center}

The exact discussion of the inertia of quadratic forms had for us only the one purpose of determining the number of real and imaginary roots of an algebraic equation by counting the positive and negative squares of the Bezoutiant. If $n$ is the degree of the equation, its Bezoutiant is a quadratic form in $n - 1$ variables, and in \S.\,73 we formed it completely for $n = 3, 4, 5$
%CONT>>
